% !PRINTED 128
\weberpage{128}
%<<CONT

or at least exceeds none of the others. By theorem 5 of \S\,38 a value $\alpha$ of $x$ can then be determined so that the absolute value $a$ of $f(\alpha)$ is smaller than $b_1$, hence
\begin{equation}
 a<b_1\le b_2\le b_3\le\cdots\le b_{n-1}.
\tag{5}
\end{equation}

We now bound a domain $G$ for the variable $x$ in such a way that \gesperrt{outside} this domain the absolute value of $f(x)$ is always \gesperrt{greater} than $a$, so that $G$ contains all points $x$ for which $|f(x)|<a$ (but also other points). By \S\,31, IV this is possible in the following way. About the origin as centre draw a circle $(R)$ of sufficiently large radius $R$ containing the points $\alpha,\beta_1,\beta_2,
\ldots,\beta_{n-1}$, and cut out of this circle the points $\beta_1,\beta_2,
\ldots,\beta_{n-1}$, at which $|f(x)|$ is certainly greater than $a$, by circular envelopes of sufficiently small but non-zero radii $\varrho_1,\varrho_2,
\ldots,\varrho_{n-1}$, so that within all these circles $(\varrho_1),(\varrho_2),\ldots,(\varrho_{n-1})$ the absolute value $|f(x)|$ remains greater than $a$ (\S\,31, V). None of these circles then encloses the point $\alpha$, for at that point $|f(x)|=a$. The connected piece of surface bounded by the circles $(R),(\varrho_1),\ldots,(\varrho_{n-1})$ is the domain $G$.

% Exact 300-dpi crop from the registered Volume-I scan, printed Fig. 3.
\weberfigureimage{weber-b1-fig03.png}{0.58}

We now determine a sequence of numbers
\[
 \alpha,\alpha',\alpha'',\alpha''',\ldots
\]
in such a way that the absolute values of $f(\alpha),f(\alpha'),f(\alpha''),\ldots$, which we denote by
\[
 a,a',a'',a''',\ldots,
\]
always decrease. The points $\alpha,\alpha',\alpha''$ so determined all remain in the interior of the domain $G$, since for them $|f(x)|<a$.

For this we use a procedure quite similar to that used for the proof of theorem 5 in \S\,38, but simplified by the fact that $f'(\alpha)$ is already different from zero. If $\delta$ denotes a positive proper fraction, still to be determined more precisely, we put in the expansion
%CONT>>
